\scanpage{200}
ALGÈBRES de LIE SEMI-SIMPLES \mspage{(5)}

Tribu n° \unsure{59|57|53}

\uline{n° actuellement 4} - Dire ``produit fini'' dans la déf. 3. Dire que $\operatorname{Autel}(\mathrm{g})$ est invariant dans $\operatorname{Aut}(\mathrm{g})$. Au milieu de la p. 12 raisonner sur $\underline{\mathrm{g}}/\underline{\mathrm{h}}$, (au lieu de Fitting).

Ajouter `` $(\mathrm{t}\in\mathrm{k})$ '' à la l.12 p. 13 ; dire ``une seule valeur'' à la l.14. Il s'agit d'une inclusion à la l.18. Dire `` image par $\mathrm{F}$ '' à la l. 19.

Dans l'énoncé du th. 2 dire ``le groupe $\operatorname{Autel}(\mathrm{G})$'' opère transitivement'' sur les Cartans. Dire ``la conclusion de (iii)'' dans la remarque p. 14.

Ajouter le cor. suivant : soient $\underline{\mathrm{h}}\subset\underline{\mathrm{g}}'\subset\underline{\mathrm{g}}$ ; si $\underline{\mathrm{g}}'$ a même rang que $\underline{\mathrm{g}}$ et contient un élément régulier de $\underline{\mathrm{g}}$, alors $\underline{\mathrm{h}}$ est une Cartan de $\underline{\mathrm{g}}$ si et seult. si c'est une Cartan de $\underline{\mathrm{g}}'$, (utiliser la prop. ajoutée au n° 3).

\uline{n° 6} - Dire partout ``commutative'', (au lieu de ``abélienne''). Ajouter ``car nilpotente'' à la l.9, p. 16. Utiliser la prop. ajoutée au n° 3 dans la prop. 14.

\uline{n° 7} - Style du th. 4 : ``le groupe des\ldots\ opère transitivement\ldots''. Idem dans le cor. p. 20. Dire ``composante neutre'' à la p. 21, l.3. Ajouter que, si $\underline{\mathrm{g}}$ est une alg. Lie sur $\mathrm{K}$, et si $\mathrm{k}$ est un sous-corps de $\mathrm{K}$, les Cartans de $\underline{\mathrm{g}}$ considérée comme algèbre sur $\mathrm{k}$ sont les mêmes que celles de $\underline{\mathrm{g}}$ considérée comme algèbre sur $\mathrm{K}$.

\S2 - \uuline{Les représentations de} $\mathrm{sl}_2$.

Eviter les points, (p. 24). Dire : ``sous-algèbre diagonale''.

\uline{n° 1} - Dans la remarque, (p. 25), dire ``(car $\mathrm{X}$ est un crochet)''. Généraliser la déf. des poids à tous les vecteurs propres de $\mathrm{H}$. Un exer. montrera les séries génératrices correspondant à (2). A la p. 25, l.2 bas, dire `` d'où, en divisant par $\mathrm{n}!$ ''.

% archive sheet number: 199
